% !TEX root = ../thesis.tex

\chapter{Introduction}\label{ch:introduction}

\gls{remotesensing} is an irreplaceable source of archaeological information, as it allows researchers to study large landscapes without relying entirely on direct fieldwork or ground observations. Among \gls{remotesensing} technologies, airborne \gls{lidar} has revolutionized archaeological exploration, particularly in areas with dense vegetation and sites where a "look under the hood" of it is required. With the technological advancements that have enabled the creation of high spatial resolution models with minimized impact of above-ground objects, \gls{lidar} data allows for the identification of complex and barely noticeable \gls{anthropogenicmicrorelief} features. Such historic modifications of the terrain as burial mounds, platforms, terraces, embankments, roads may be difficult or impossible to detect using traditional optical imagery~\cite{thompson_detecting_2020,kokalj_standardizing_2025}.

However, the increasing accessibility to \gls{lidar} data has also uncovered certain fundamental methodological obstacles, as the number of topographical datasets being acquired is growing faster than expert's capacity to manually review, annotate, and verify them. Despite the continuous development of complex analysis methods, traditional archaeological interpretation still relies significantly on expert visual analysis of the obtained relief data, and this role remains critically important. The disadvantage of this is that manual interpretation is labor-intensive, difficult to scale to wide areas, and also remains vulnerable to differences in experience, selection of visualization, annotation conventions, and localized archaeological and historical context~\cite {trier_using_2019,verschoof-van_der_vaart_combining_2020,kokalj_standardizing_2025}. Precisely for this reason, one of the main challenges is not to completely exclude the expert from the workflow, but rather the development of computational approaches that can support or oppose the expert's judgement by focusing on specific areas of interest to reduce routine brute-force effort in searching for potential features while ensuring a more systematic and reproducible data processing approach.

Deep learning architectures, in particular convolutional neural networks, offer a promising framework to address this challenge. The use of \gls{lidar} data in the context of archaeology fits seamlessly with neural networks, and the ability to use almost any two-dimensional dataset for the purpose of training allows for the usage of all currently available types of terrain data representations. While the tasks addressed by these models differ in terms of type, the result in all cases leads to the identification of anthropogenic objects. Recent studies have already demonstrated the potential of approaches utilizing model families such as U-Net, \gls{rcnn}, \gls{yolo}, and other \gls{cnn} architectures for discovering archaeological objects in data obtained using \gls{lidar}~\cite{bundzel_semantic_2020,somrak_learning_2020,guyot_objective_2021,zhang_unveiling_2024}. The main advantage of such methods lies in accelerating the processing of a large study area, in contrast to manually reviewing the scanned territory, moreover, they help to highlight individual zones with a higher probability of potential archaeological objects.

Despite the existing potential and the promising results already reported, the applicability of deep learning algorithms for the processing of archaeological \gls{lidar} data remains methodologically immature. Published studies often differ in their approaches to key steps, such as data preparation, \gls{terrainvisualization}, label construction, model architecture, sampling strategies, \gls{lossfunction}, threshold selection, and evaluation procedures. As a consequence of the lack of established protocols, it is difficult to determine which specific components of the workflow have a significant impact on the overall quality of predictions, and which are merely secondary aspects of implementation. The lack of standardization directly limits the comparability of studies and complicates the transfer of methods across different regions, site types, and data resolutions, forcing authors to manually test each individual component to determine its impact on achieving an acceptable result~\cite{fiorucci_deep_2022,kokalj_standardizing_2025}.

Several specific problems are especially relevant for this thesis:
\begin{enumerate}
    \item \textbf{Choice of \gls{inputrepresentation}.} \gls{lidarderiveddata} terrain data can be represented through different raster products and visualizations, such as \gls{dem}, \gls{hillshade}, \gls{slope}, \gls{skyviewfactor}, openness, or combined multi-band visualizations. Each representation emphasizes different aspects of the terrain surface. Consequently, the selected representation can significantly influence which archaeological structures and their characteristics are visible to the model and how well they are separated from the surrounding natural terrain~\cite{thompson_detecting_2020,kokalj_standardizing_2025}.

    \item \textbf{Quality and form of reference labels.} Archaeological annotations are usually produced by experts for interpretation and documentation purposes, not primarily for training machine learning models. Therefore, the label geometry may be incomplete, conservative, spatially imprecise, or inconsistent with the exact topographic expression visible in selected terrain representation. For segmentation models, these properties directly affect the training signal and may lead to noisy or ambiguous supervision~\cite{bundzel_semantic_2020,rahaman_effects_2022,burgert_effects_2022}.

    \item \textbf{Extreme \gls{classimbalance}.} Archaeological objects usually occupy only a very small fraction of all pixels in a \gls{lidar} raster. The \gls{backgroundclass} is therefore dominant, while the target class may be sparse, narrow, fragmented, or locally discontinuous. Without an appropriate sampling strategy and \gls{lossfunction}, the model can learn to prefer the \gls{backgroundclass} and suppress \gls{archaeologicalfeature}~\cite{diakogiannis_resunet-_2020,zhou_dynamic_2023}.

    \item \textbf{Evaluation of practical usefulness.} Standard \gls{pixelwisemetric}, such as \gls{intersectionoverunion}, \gls{precision}, \gls{recall}, or \gls{fscore}, are necessary for quantitative comparison, but they do not always fully reflect the archaeological value of predictions. A segmentation result that is slightly shifted from the reference label may receive a low pixel-wise score, even though it can still be useful as a candidate for expert inspection or field verification. For this reason, quantitative evaluation should be complemented by qualitative and \gls{gis}-based assessment~\cite{fiorucci_deep_2022,guyot_objective_2021}.
\end{enumerate}

\section{Motivation}

The motivation of this thesis follows from the gap between the potential of deep learning methods and the current methodological uncertainty in their application to archaeological \gls{lidar} data. Existing research shows that neural networks can support archaeological prospection, but it also shows that the final result depends on many decisions made before, during, and after model training. These decisions include the choice of \gls{inputrepresentation}, the transformation of expert annotations into training masks, the sampling of highly imbalanced data, the selection of a \gls{lossfunction}, the optimization of the prediction threshold, and the final interpretation of the output in a \gls{gis} environment.

In practical terms, the problem is that a trained segmentation model is not useful only because it produces a numerical score. It is useful only if its predictions can help an archaeologist inspect large areas more efficiently, identify candidate features, and decide where further verification may be needed. Therefore, the main motivation is to build and evaluate a reproducible experimental workflow that makes these decisions explicit and comparable. Such a workflow can clarify which components improve segmentation quality, which configurations are unstable, and how the same methodology behaves in different archaeological and geographical contexts.

This thesis is motivated by two complementary data domains: Maya archaeological structures in Guatemala and archaeological terrace systems in Slovakia. These datasets differ in object morphology, spatial resolution, annotation type, and landscape context. Their comparison makes it possible to study not only whether \gls{semanticsegmentation} can be applied to archaeological \gls{lidar} data, but also how sensitive the workflow is to changes in the data domain. This is important because a method that works only in one narrowly defined case has limited scientific and practical value.

\section{Formulation of the problem}

This thesis addresses the design of a reproducible methodology and conducts a comparative analysis of the components of the \gls{semanticsegmentation} workflow for anthropogenic \gls{archaeologicalfeature} in terrain data acquired using airborne \gls{lidar}. Furthermore, the study centers on segmentation in two archaeological contexts: structures related to Mesoamerican culture in a Guatemalan dataset, and micro-relief landforms that resemble a series of terraces in a Slovakian dataset.

The objective is defined as a \gls{binarysemanticsegmentation} task, in which each valid pixel in a two-dimensional input matrix—representing raster or elevation data—is assigned to either an \gls{archaeologicalfeature} or the background. The methodology for studying this task involves several interconnected steps: preparation of input data obtained via \gls{lidar}; conversion of existing expert annotations into training masks; construction of a \gls{tile}-based dataset suitable for deep learning; training of segmentation models under controlled experimental conditions; generating \gls{thresholdbasedinference} \gls{binaryprediction}; and evaluating results using both quantitative metrics and qualitative \gls{gis}-based assessments.

Accordingly, the research questions addressed in this thesis do not intend to design a fully self-sufficient system for scanning landscapes for potential anthropogenic structures that would replace expert-based methods of interpretation. Instead, it is grounded in the design, implementation, and evaluation of a controlled workflow that directly utilizes deep learning algorithms to support decision-making in archaeological exploration by producing reproducible segmentation results and determining the impact of individual methodological components on the final outcome. That is, the expected outcome is an experimentally validated assessment of how different data representations, label formats, sampling strategies, \gls{lossfunction}, and evaluation procedures affect the capacity of the model to produce meaningful predictions of potential archaeological occurrences present in topographic data obtained using airborne \gls{lidar}.
